===============================================================================
\item メタノール(\ce{CH3OH}，32.04\ g/mol)\cite{d}\\
    \hspace{0pt}・常温常圧下では，無色透明の液体であり，臭気を持つ。\\
    \hspace{0pt}・水やエタノールに極めて溶けやすい。\\
    \hspace{0pt}・酸化剤と反応し，火災や爆発の危険をもたらす。\\
    \hspace{0pt}・第一級アルコールである。\\

\begin{minipage}[b]{0.30\textwidth}
    \centering
    \chemfig[atom sep=2.5em]{-[,0.8]OH}
    \caption{メタノール}
    \label{fig:CH3OH}
\end{minipage}  

\bibitem{d}  厚生労働省：メタノール，職場のあんぜんサイト(オンライン)，\\
             入手先\url{<https://anzeninfo.mhlw.go.jp/anzen/gmsds/67-56-1.html>}
             (参照2025-4-18).
===============================================================================
\item ヘキサン(\ce{C6H14}，86.18\ g/mol)\cite{d}\\
    \hspace{0pt}・常温常圧下では，無色透明の液体であり，臭気を持つ。\\
    \hspace{0pt}・エタノールに極めて溶けやすいが，水には溶けにくい。\\
    \hspace{0pt}・酸化剤やハロゲンと反応し，火災や爆発の危険をもたらす。\\
    \hspace{0pt}・引火点が極めて低い引火性の物質であり，火災の危険をもたらす。\\

\begin{minipage}[b]{0.30\textwidth}
    \centering
    \chemfig[atom sep=2.5em]{-[:-30]-[:30]-[:-30]-[:30]-[:-30]}
    \caption{ヘキサン}
    \label{fig:C6H14}
\end{minipage}

\bibitem{d}  厚生労働省：ヘキサン，職場のあんぜんサイト(オンライン)，\\
             入手先\url{<https://anzeninfo.mhlw.go.jp/anzen/gmsds/110-54-3.html>}
             (参照2025-4-18).
===============================================================================
\item トルエン(\ce{C7H8}，92.14\ g/mol)\cite{d}\\
    \hspace{0pt}・常温常圧下では，無色透明の液体であり，臭気を持つ。\\
    \hspace{0pt}・エタノールやエーテルに極めて溶けやすいが，水には溶けにくい。\\
    \hspace{0pt}・強酸化剤と激しく反応し、火災や爆発の危険をもたらす。\\
    \hspace{0pt}・引火点が極めて低い引火性の物質であり，火災の危険をもたらす。\\
    \hspace{0pt}・ベンゼン環を持つ有機化合物である。\\

\begin{minipage}[b]{0.30\textwidth}
    \centering
    \chemfig[atom sep=2.5em]{*6(-=-(-[,0.8])=-=)}
    \caption{トルエン}
    \label{fig:C7H8}
\end{minipage}

\bibitem{d}  厚生労働省：トルエン，職場のあんぜんサイト(オンライン)，\\
             入手先\url{<https://anzeninfo.mhlw.go.jp/anzen/gmsds/0045.html>}
             (参照2025-4-18).
===============================================================================
\item サリチル酸(\ce{C7H6O3}，\ g/mol)\cite{d}\\
    \hspace{0pt}・常温常圧下では，白色の結晶〜結晶性粉末である。\\
    \hspace{0pt}・水に僅かに溶け，熱湯やエタノール，エーテルに溶けやすい。\\
    \hspace{0pt}・光に敏感であり，日光の下で徐々に変色する。\\
    \hspace{0pt}・水溶液は弱酸性の物質である。\\

\begin{minipage}[b]{0.40\textwidth}
    \centering
    \chemfig[atom sep=2.5em]{*6((-H@{x6}O)-=-=-(-C(=[:90]O)(-[:210]HO))=)}
    \caption{サリチル酸}
    \label{fig:C7H6O3}
\end{minipage}

\bibitem{d}  厚生労働省:サリチル酸，職場のあんぜんサイト(オンライン)，\\
             入手先\url{<https://anzeninfo.mhlw.go.jp/anzen/gmsds/69-72-7.html>}
             (参照2025-5-8).
===============================================================================
\item 炭酸ナトリウム(\ce{Na2CO3}，105.99\ g/mol)\cite{d}\\
    \hspace{0pt}・常温常圧下では，白色の粉末であり，吸湿性がある。\\
    \hspace{0pt}・水に溶けやすく，エタノールにほとんど溶けない。\\
    \hspace{0pt}・マグネシウム，五酸化リンと反応し，爆発の危険をもたらす。\\
    \hspace{0pt}・水溶液は強塩基性の物質である。\\

\bibitem{d}  厚生労働省:炭酸ナトリウム，職場のあんぜんサイト(オンライン)，\\
             入手先\url{<https://anzeninfo.mhlw.go.jp/anzen/gmsds/497-19-8.html>}
             (参照2025-5-8).
===============================================================================
\item 無水酢酸(\ce{(CH3CO)2O}，102.09\ g/mol)\cite{d}\\
    \hspace{0pt}・常温常圧下では，無色透明の液体であり，臭気を持つ。\\
    \hspace{0pt}・エーテルに極めて溶けやすく，エタノールと水に反応して溶ける。\\
    \hspace{0pt}・水もしくは水分が存在すると強い腐食性を示す。\\
    \hspace{0pt}・強酸化剤の物質や強塩基性の物質と激しく反応する。\\

\begin{minipage}[b]{0.40\textwidth}
    \centering
    \chemfig[atom sep=2.5em]{-[:30](=[:90]O)-[:-30]O-[:30](=[:90]@{x1}O)-[:-30]}
    \caption{無水酢酸}
    \label{fig:(CH3CO)2O}
\end{minipage}

\bibitem{d}  厚生労働省:無水酢酸，職場のあんぜんサイト(オンライン)，\\
             入手先\url{<https://anzeninfo.mhlw.go.jp/anzen/gmsds/0965.html>}
             (参照2025-5-11).
===============================================================================
\item アニリン(\ce{C6H5NH2}，93.13\ g/mol)\cite{d}\\
    \hspace{0pt}・常温常圧下では，無色～淡黄色の液体であり，特異臭がある。\\
    \hspace{0pt}・エタノール及びエーテルに極めて混じりやすく，水にやや混じりやすい。\\
    \hspace{0pt}・光に暴露して茶色に変色する。
    \hspace{0pt}・可燃性があり，火災や爆発の危険をもたらす。\\
    \hspace{0pt}・強酸や強酸化剤と反応する。\\

\begin{minipage}[b]{0.30\textwidth}
    \centering
    \chemfig[atom sep=2.5em]{*6(-=-(-[,0.8]NH2)=-=)}
    \caption{アニリン}
    \label{fig:C6H5NH2}
\end{minipage}

\bibitem{d}  厚生労働省:アニリン，職場のあんぜんサイト(オンライン)，\\
             入手先\url{<https://anzeninfo.mhlw.go.jp/anzen/gmsds/62-53-3.html>}
             (参照2025-5-18).
===============================================================================
\item 酢酸ナトリウム三水和物(\ce{CH3COONa.3H2O}，136.08\ g/mol)\cite{d}\\
    \hspace{0pt}・常温常圧下では，白色の結晶〜結晶性粉末である。\\
    \hspace{0pt}・水に極めて溶けやすく，エタノールにやや溶けやすい。\\
    \hspace{0pt}・強酸化剤と激しく反応する。\\

\bibitem{d}  林純薬工業:酢酸ナトリウム三水和物，試薬ダイレクト(オンライン)，\\
             入手先\url{<https://direct.hpc-j.co.jp/shop/g/g16605-1/>}
             (参照2025-5-18).
===============================================================================




===============================================================================




===============================================================================



===============================================================================


===============================================================================




===============================================================================
